\section{Motivation and Challenges}
\label{sec:motivation}

The resulting management problem is to decide which KV to retain and
in which tier. For a multi-agent service, the benefit of these decisions
depends on how retained state supports both normal execution and workflow
completion after a failure. Existing fault-tolerant serving and
workflow-aware KV management provide useful mechanisms, but their
objectives leave this joint decision unresolved.

We measure the application-level outcome by \emph{workflow completion
 time} (WCT), $\mathrm{WCT}_w=c_w-a_w$, where $a_w$ is workflow arrival
 and $c_w$ is the time all required outputs become ready. This includes
 model and tool execution, queueing, and dependency waiting.

\subsection{Limitations of Existing Approaches}
\label{sec:motivation-limitations}

\begin{figure*}[t]
  \centering
  \captionsetup[subfigure]{font=normalsize,labelfont=normalfont,
    textfont=normalfont,justification=centering,singlelinecheck=false}
  \begin{subfigure}[t]{2.48in}
    \vspace{0pt}
    \begin{minipage}[c][1.65in][c]{\linewidth}
      \centering
      \includegraphics[width=\linewidth]{figures/workflow.pdf}
    \end{minipage}
    \caption{Four-branch workflow}
    \label{fig:motivation-workflow}
  \end{subfigure}\hfill
  \begin{subfigure}[t]{1.86in}
    \vspace{0pt}
    \begin{minipage}[c][1.65in][c]{\linewidth}
      \centering
      \includegraphics[width=\linewidth]{figures/workflow-gaps.pdf}
    \end{minipage}
    \caption{Workflow-maximum gaps}
    \label{fig:motivation-branch-gap}
  \end{subfigure}\hfill
  \begin{subfigure}[t]{2.52in}
    \vspace{0pt}
    \begin{minipage}[c][1.65in][c]{\linewidth}
      \centering
      \includegraphics[width=\linewidth]{figures/recovery-exposure.pdf}
    \end{minipage}
    \caption{Recovery cost and WCT}
    \label{fig:motivation-recovery-exposure}
  \end{subfigure}
  \caption{Workflow timing and conditional recovery impact.
    (a) Four trace-derived parallel agents A1-A4 with schematic Main agents A0/A5.
    Bars preserve timing; dashed tails assume an all-branch join.
    (b) Empirical CDF of workflow-level gaps over 121 workflows. For each parallel group with temporally inferred fan-in, we subtract the shortest agent execution time from the longest, then retain the largest gap per workflow. 
   (c) Each horizontal bar represents one agent (A1--A4) and shows its
extra delay to reconstruct KV after a failure, measured by full
prefill with Qwen2.5-7B-Instruct-1M. The hatched portion is absorbed
by slack; the solid portion is the calculated WCT increase.}

  \label{fig:motivation}
\end{figure*}

\paragraph{Limitation 1: Fault tolerance without workflow-aware priorities.}
Fault-tolerant serving systems preserve state and coordinate recovery
for interrupted requests. D\'ej\`aVu streams KV checkpoints, and GhostServe
uses host-memory parity to reduce reconstruction
work~\cite{strati2024dejavu,jayakody2026ghostserve}.
LUMEN further coordinates checkpoint placement, recovery routing and
capacity restoration according to worker load~\cite{cao2026lumen}.
These mechanisms address state loss and recovery contention, but do not
use multi-agent dependencies to determine which requests' faster
recovery would advance workflow completion. Request recovery cost and
worker load alone cannot express this benefit.

\label{sec:motivation-workflow}
Figure~\ref{fig:motivation}(\subref{fig:motivation-workflow}) illustrates
why workflow position matters. A2 finishes last, 151.6\,s after the fork,
whereas A1 finishes at 60.4\,s. If the following Main stage requires all
four outputs, A1 can absorb 91.2\,s of additional delay before postponing
the join, whereas any delay on A2 postpones it. This example comes from the Claude Code traces~\cite{semianalysis2026cctraces} and has the fork--join structure illustrated in Figure~\ref{fig:serving_model}. We call the additional
delay a branch can absorb without postponing the required output its
\emph{slack}. 

This situation where agents with parallel execution have different completion times is very common in multi-agent applications, and the time differences can be significant. We analyze the Claude code trace~\cite{semianalysis2026cctraces} related information and visualize this time difference in Figure~\ref{fig:motivation}(\subref{fig:motivation-branch-gap}).
This figure reports each workflow's maximum longest-minus-shortest branch-duration gap,
with a median of 138\,s and a 90th percentile of 1,173\,s. When the Main agent waits for all required branches, differences in branch completion times can leave shorter branches with slack. Additional delays within this slack do not postpone the join or
increase WCT.

These differences in slack affect how KV reconstruction delays workflow
completion after a failure.
Figure~\ref{fig:motivation}(\subref{fig:motivation-recovery-exposure})
shows reconstruction times measured by full prefill with
Qwen2.5-7B-Instruct-1M using sparse
attention~\cite{qwen2025qwen25oneM,jiang2024minference},
and their calculated impact on WCT.
We consider each agent's interruption separately, with other execution
unchanged and the Main agent waiting for all four branches.
A1's 33.2\,s reconstruction delay is fully absorbed by its 91.2\,s slack,
leaving WCT unchanged. In contrast, A2 has no slack, so its 32.6\,s
reconstruction delay increases WCT by the same amount.
Thus, similar reconstruction costs can have very different effects
on workflow completion. Recovery priorities should therefore account
for a request's position in the workflow, in addition to its
reconstruction cost.

\paragraph{Limitation 2: Workflow-aware KV management ignores failures.}
Workflow-aware serving systems already exploit the dependencies
highlighted above, but to accelerate normal execution. One line of work
hands workflow structure to the scheduler: Parrot exposes inter-request
dependencies for application-level scheduling~\cite{lin2024parrot}, Ayo
orchestrates applications as primitive-level dataflow
graphs~\cite{tan2025ayo}, Autellix and SAGA schedule whole agent programs
rather than individual calls~\cite{luo2025autellix,guo2026saga}, and
Cortex and Murakkab pool and configure resources from the same
structure~\cite{pagonas2025cortex,chaudhry2026murakkab}. Another line
manages KV with it directly: InferCept and Continuum decide whether KV is
kept, swapped or discarded during a tool call, and for how
long~\cite{abhyankar2024infercept,li2025continuum}; TokenCake offloads KV
during tool stalls, prefetches it before the next call, and reserves GPU
memory for critical-path agents~\cite{bian2025tokencake}; KVFlow,
CacheScout and PBKV order eviction and prefetching by predicted upcoming
agent invocations~\cite{pan2025kvflow,zhang2026cachescout,zheng2026pbkv};
Pythia predicts future steps and prompt contents to stage KV ahead of
use~\cite{yu2026pythia}; and TOPAS jointly selects resident prefixes and
admitted requests from remaining-path progress, reuse, and movement or
preemption costs~\cite{ni2026topas}. These policies value retained KV by
the requests expected to read it next, and assume that the worker holding
it keeps running.

\begin{figure}[t]
  \centering
  \includegraphics[width=\columnwidth]{figures/gpu-server-availability.pdf}
  \caption{GPU-server availability reconstructed from InfiniteHBD's
    production-training trace~\cite{shou2025infinitehbd}: 400 servers minus
    those within recorded outage intervals. Records include GPU, network,
    power/cooling, and software faults; unknown causes, tests,
    troubleshooting, and changes are also retained. Overlapping intervals
    per server are merged. Availability is measured at the node level.}
  \label{fig:gpu-server-availability}
\end{figure}

Yet failures are not rare, as Figure~\ref{fig:gpu-server-availability}
shows for a production training cluster: 584 recorded unavailability
events affect 231 of 400 servers in about 349 days~\cite{shou2025infinitehbd}. When a server fails, the agent requests running on it are
interrupted together with their GPU-resident KV and any copies swapped
to that host's
memory~\cite{abhyankar2024infercept,pan2025kvflow,bian2025tokencake}.
What remains on other workers or in shared storage are the fixed prompts
and shared prefixes selected by their next expected
read~\cite{pan2025kvflow,yu2026pythia,zhang2026cachescout}; the
interrupted request's history up to the failure is not among them. The
Claude Code trace bounds how much such surviving prefixes could help:
same-agent prefixes cover 97.9\% of all request input (98.7\% for
main-agent and 93.9\% for subagent requests), whereas prefixes from other
agents cover at most 10.1\% (4.4\% and 36.2\%), largely overlapping the
former (Figure~\ref{fig:agent-prefix-reuse}). Most of an interrupted request's
input is thus its own agent's history, which resided with that agent on
the failed worker; it must be reconstructed from the prompt and token
log, at the costs of Figures~\ref{fig:motivation-recompute}
and~\ref{fig:motivation-recovery-exposure}, and which agent it lands on
decides whether that time is absorbed by slack or added to WCT. The
root cause is that failure is absent from these systems' model: whether
a KV object is worth keeping, where, and for how long is decided on the
premise that its worker does not fail. Swapping to host memory frees GPU
capacity rather than escaping a failure, shared storage holds completed
contexts for later reuse rather than a running request's history, and a
copy that would be read only after a failure is worth nothing in their
retention scores. So although these systems already hold the information
needed to tell which agent cannot afford this reconstruction---the
dependency graphs, next-call predictions and remaining paths they extract
are exactly what slack is derived from, and TokenCake even uses them to
reserve GPU memory for critical-path agents~\cite{bian2025tokencake}---that
information is never used to decide, before a failure, for which agent,
at which tier, and for how long a resumable copy should be kept.

\begin{figure}[t]
 \centering
 \includegraphics[width=\columnwidth]{figures/agent-prefix-reuse.pdf}
 \caption{Token-weighted potential prefix coverage in the Claude Code
 trace~\cite{semianalysis2026cctraces}: all 98,827 requests, split into
 56,798 main-agent and 42,029 subagent requests. For each request,
 coverage is the longest input prefix shared with an earlier completed
 request of the same workflow and model, summed over tokens. Same- and
 other-agent bars overlap: other agents add only 0.37\%, 0.01\% and
 2.05\% beyond the best same-agent prefix. These are reuse opportunities,
 not runtime cache hits.}
 \label{fig:agent-prefix-reuse}
 \Description{Three groups show all requests, main-agent requests and
 subagent requests. Same-agent coverage is 97.9, 98.7 and 93.9 percent;
 other-agent coverage is 10.1, 4.4 and 36.2 percent, respectively.}
\end{figure}

\label{sec:motivation-long-context}
The two limitations are mirror images. Fault-tolerant serving decides
\emph{how} to preserve and restore state, but ranks interrupted requests
by recovery cost and worker load. Workflow-aware caching decides
\emph{which} state a workflow will read next, but assumes the worker
holding it survives. Neither answers the question a multi-agent service
faces when a worker fails: which agents' histories must be recoverable
quickly for their workflows to complete on time, and how much of the
finite HBM, DRAM and storage that protection may take from normal reuse.
Combining the two is not a matter of adding slack to a recovery router
or a checkpoint path to a cache. Both objectives draw on the same tiers,
the state each values differs, and both change while agents run.
AegisKV takes this joint decision as its problem. It assumes checkpoint
transport and tiered stores as in the systems above, and contributes the
placement policy---which KV objects each tier holds, on whose behalf,
and when they move---together with the bounded controller and safe
transitions that keep that policy correct online. Realizing it raises
three challenges.

\subsection{Design Challenges}
\label{sec:motivation-challenges}

\paragraph{Challenge 1: Slack requires predicting each agent's remaining execution time online.}
Slack is a statement about the future: how much longer an agent's
branch may run before it, rather than its slowest sibling, gates the
join. The examples in Section~\ref{sec:motivation-limitations} read it
off completed traces. Online, a protection decision needs the remaining
execution time of every running agent at the moment it is made---how
many model calls and tool waits are left, and how long each will
take---and none of this is observed. Tool durations are outside the
serving system's control, the number of further turns depends on outputs
not yet generated, and the branch that gates the join changes as
siblings finish, so one agent's slack shifts whenever another agent
progresses. The quantity to be predicted also spans a wide range:
workflow-maximum gaps run from seconds to over a thousand seconds
(median 138\,s, 90th percentile 1,173\,s in
Figure~\ref{fig:motivation}(\subref{fig:motivation-branch-gap})), so no
fixed rule separates urgent from tolerant branches, and neither recency
nor reconstruction cost stands in for the prediction: A1 and A2
reconstruct in similar time, yet only A2 delays the join. The system
must therefore estimate remaining time from causal features of partial
execution alone, refresh the estimate as agents advance, and derive
priorities that remain useful under prediction error and fall back
safely when the estimate is invalid.

\begin{figure}[t]
  \centering
  \includegraphics[width=\columnwidth]{figures/recovery-cost.pdf}
  \caption{Measured KV reconstruction and reuse costs for
    Qwen2.5-7B-Instruct-1M on two A100 PCIe 40\,GB GPUs
    (TP2, BF16, batch one), with sparse attention above 32K.
    Full prefill starts from empty KV; reuse starts from $(C-1)$ cached
    tokens and computes one. KV reconstruction time includes any
    retrieval and computation until all $C$ tokens' KV is ready.}
  \label{fig:motivation-recompute}
  \Description{Four measured KV reconstruction-time curves for Qwen2.5-7B-Instruct-1M
    from 8K to 1,010,000 tokens: blue circles for resident HBM reuse,
    orange squares for DRAM reuse, purple diamonds for shared-storage
    reuse, and red triangles for full prefill. Both axes are logarithmic.
    At the longest context, the four paths take 1.305, 2.663, 10.156 and
    571.21 seconds, respectively.}
\end{figure}

\paragraph{Challenge 2: KV priority must combine reuse value with recovery benefit, then decide placement.}
Existing systems rank KV by one criterion: workflow-aware caching by the
requests expected to read it next, fault-tolerant serving by the cost of
reconstructing it (Section~\ref{sec:motivation-limitations}). Neither is
the value a multi-agent service needs. A KV object is worth keeping for
two reasons that coexist and can disagree. Normal reuse saves prefill on
the agent's next call whether or not its worker fails. Recovery saves
reconstruction only if a failure occurs, and matters for WCT only to the
extent that the saved time exceeds the agent's slack from
Challenge~1. High slack must therefore not delete history that later
normal calls would recompute, and a low reconstruction cost must not
demote the one agent whose reconstruction would land on the critical
path. Recovery benefits also do not add up: one worker failure
interrupts several co-located agents whose recoveries share readers,
links and target GPUs, and accelerating one of two late branches leaves
their join unchanged. The redefined priority then has to decide storage
management. Retaining a history in any tier removes reconstruction: at
1M tokens, 571.21\,s becomes 10.156\,s from shared storage
(Figure~\ref{fig:motivation-recompute}). A faster tier only shortens
retrieval, to 2.663\,s in DRAM and 1.305\,s in HBM, at a very different
price: that history occupies 57.92\,GB, over 70\% of the two A100s'
combined HBM, and because interrupted agents' histories are largely
unique (Figure~\ref{fig:agent-prefix-reuse}), sharing does not shrink
this cost. A history need not move as a unit, however. Long-context
models increasingly select blocks, tokens or compressed entries per
query (Figure~\ref{fig:sparse-releases}), and their accesses
concentrate: in NOSA-8B, NSA-8B and MiniCPM4.1-8B, about 20\% of
historical blocks receive 87.64\%, 74.76\% and 56.54\% of accesses
(Figure~\ref{fig:design-insights}). A small hot subset in a fast tier
could serve most reads while the complete history stays in cheaper
tiers, provided the subset is chosen from causal history alone and a
missed block remains retrievable. Placement must therefore turn the
combined priority into decisions per tier and per block---which objects
each tier holds, on whose behalf and for how long---separating
retention that avoids reconstruction from acceleration that shortens
retrieval, within finite tier budgets.

\begin{figure}[t]
 \centering
 \includegraphics[width=\columnwidth]{figures/sparse-models.pdf}
 \caption{Selected sparse/hybrid model releases and reported context
 lengths. IDs: 1 MiniCPM4~\cite{minicpm4model};
 2 MiniCPM4.1~\cite{openbmb2025minicpm41};
 3 DeepSeek-V3.2-Exp~\cite{deepseek2025v32exp};
 4 DeepSeek-V3.2~\cite{deepseek2025v32model};
 5 GLM-5~\cite{zai2026glm5model};
 6 MiniCPM-SALA~\cite{openbmb2026sala};
 7 DeepSeek-V4 preview~\cite{deepseek2026v4model};
 8 MiniMax-M3~\cite{minimax2026m3};
 9 GLM-5.2~\cite{zai2026glm52};
 10 LongCat-Flash-Lite-Sparse~\cite{longcat2026sparse}.
 Dates are announcements, or first official repository dates for
 IDs 5, 9 and 10. Markers identify the sparse branch's selection unit,
 including SALA's block branch and V4's compressed sparse branch;
 they do not imply uniform attention across layers. This is a selected
 survey, not an adoption-share or sparsity-rate estimate.}
 \label{fig:sparse-releases}
 \Description{Ten selected sparse or hybrid attention releases between
 June 2025 and July 2026, with reported contexts from 32K to 1M.
 Circles, squares and diamonds distinguish block, token and compressed
 entry selection. This is a literature survey, not a sparsity-rate curve.}
\end{figure}

\begin{figure}[t]
 \centering
 \includegraphics[width=\columnwidth]{figures/kv-access-concentration.pdf}
 \caption{Native KV access concentration for NOSA-8B~\cite{huang2026nosav2},
 NSA-8B (MiniCPM-4-8B-NSA)~\cite{fasa2026minicpm4nsa}, and
 MiniCPM4.1-8B~\cite{openbmb2025minicpm41}. We use the final 512 queries
 of matched 32,704-token prefixes from three PG19 texts. Each query
 counts a historical 64-token block once across heads. Rank the 503
 blocks within each layer/window; average access shares over all
 32 layers, then texts. Bands show text min--max; the diagonal denotes
 uniform access. Compressed/selector state is excluded. This is
 retrospective concentration, not prediction accuracy or recovery speedup.}
 \label{fig:design-insights}
 \Description{Cumulative historical block access curves for NOSA-8B,
 NSA-8B and MiniCPM4.1-8B, averaged over every layer and three texts.
 About 20 percent of blocks account for 87.64, 74.76 and 56.54 percent
 of accesses, respectively; all three curves exceed the uniform baseline.}
\end{figure}

\paragraph{Challenge 3: Extra protection must not disrupt normal inference.}
Every copy kept for recovery beyond what normal serving would keep is
extra work on the serving path. Writing checkpoints and moving them
between tiers consumes the DRAM and storage bandwidth that normal reuse
and concurrent recoveries read from, and a copy held in HBM occupies
memory that inference batches need: every HBM byte spent on protection
is taken from serving. Deciding what to protect is not free either.
Every generated token extends the unsaved tail, every tool return and
new request shifts slack, and every batch changes the HBM inference
needs, so re-scoring all resident objects on each event is prohibitive
at block granularity, while stale scores keep the wrong state in the
fast tiers. Protection must also not undermine itself: transfers are
asynchronous and the failure they guard against can strike midway;
releasing a DRAM copy before its storage replacement commits, or
counting an object as local after its owner failed, silently leaves a
checkpoint without a recoverable copy. The controller therefore needs
bounded work per event, a protection capacity that yields to inference
when foreground demand rises and grows back when it falls, migration
paced within the bandwidth budget, and transitions that keep every
committed checkpoint with a reachable owner.

\label{sec:motivation-goals}
These challenges organize the design. Workflow Tracker predicts remaining
execution time and derives slack from it for Challenge~1
(Section~\ref{sec:workflow-protection}); the slack-aware request
priority, block-level demand and renewable per-agent credit define the
combined priority and its placement for Challenge~2
(Sections~\ref{sec:placement-value}--\ref{sec:sparse-placement}); and
the lazy event loop, the AIMD Adapter for Peer capacity and versioned
transfers keep protection off the inference path for Challenge~3
(Sections~\ref{sec:placement-plan} and~\ref{sec:adapter}).
